We compare our DialogGSR with existing knowledge{\textendash}grounded dialog generation models on OpenDialKG dataset. 
Table~\ref{tab:opendialkg} shows that DialogGSR achieves the best performance in all metrics~(BLEU, ROUGE, KQA, and F1 score).
In particular, DialogGSR outperforms other baselines on KQA metrics by a large margin (4.61 on EM metric), which indicates that the proposed method generates more factually correct responses with relevant knowledge.
In addition, our method achieves a 1.53 performance gain on BLEU-1 metric compared to the best baseline method, which is an 8.61\% improvement.  
The performance gain of DialogGSR compared to SURGE, which retrieves the subgraph with a bi-encoder and uses graph neural networks for graph representations, indicates that our generative retrieval is effective in retrieving relevant knowledge and generating more accurate responses based on the retrieved knowledge. 

We also conduct experiments on KOMODIS~\cite{galetzka2020corpus} dataset.
Similar to the OpenDialKG result, Table~\ref{tab:komodis} demonstrates that our DialogGSR achieves the best performance compared to all the previous approaches.
To further validate the effectiveness of our generative subgraph retrieval, we compare the retrieval performance by path@k metrics.
Table~\ref{tab:retrieval} shows that DialogGSR achieves the best performance compared to the other baselines.  
This result indicates that our generative subgraph retrieval successfully retrieves context-relevant subgraphs from the knowledge graph by fully utilizing the power of pretrained language models. 

\subsection{Human Evaluation}
\label{sup_sec:human}
We conduct human evaluation to assess the generated responses of our dialog generation model.
The detailed process of human evaluation is in Appendix~\ref{sec:human_evaluation_detail}.
Table~\ref{tab:humaneval} shows the experimental results of the human evaluation, where DialogGSR outperforms SURGE in all the metrics (Consistency, Informativeness, Fluency). 
In particular, on the Consistency and Informativeness metrics, DialogGSR achieves statistically significant performance gains of 0.16 and 0.47 over SURGE (based on $t$-test with $p$-value $<0.05$), which indicates that our generative subgraph retrieval performs significantly better in retrieving informative knowledge compared to existing retrieval methods.
Our DialogGSR provides a relatively small performance gain of 0.11 on the Fluency metric.
Since the Fluency metric is more influenced by the language model's performance than the knowledge retrieval performance, it is reasonable to expect similar fluency scores when using the same base language model (T5-small) for fair comparisons.